% TOP_PROBLEM: {"id": 7200021, "title": "Zariski multiplicity conjecture on the topological equisingularity and equimultiplicity of varieties at singular points", "category_id": 5, "category": {"id": 5, "name": "algebraic_geometry", "display_name": "Algebraic Geometry", "description": "Geometric objects defined by polynomial equations.", "slug": "algebraic-geometry", "order_index": 5, "created_at": "2026-07-31T15:26:25.671Z"}, "status": "open"}
% FIRST_POSTED: null
% ENTRY_KIND: statement_only
% !TeX program = lualatex
\documentclass[11pt]{article}
\usepackage[margin=25mm]{geometry}
\usepackage{fontspec}
\setmainfont{Latin Modern Roman}
\usepackage{amsmath,amssymb,mathtools}
\usepackage{unicode-math}
\setmathfont{Latin Modern Math}
\usepackage{xurl,hyperref}
\hypersetup{colorlinks=true,urlcolor=blue}
\setcounter{secnumdepth}{0}
\setlength{\parindent}{0pt}
\setlength{\parskip}{5pt}
\setlength{\emergencystretch}{3em}
\newcommand{\R}{\mathbb R}
\newcommand{\N}{\mathbb N}
\newcommand{\Z}{\mathbb Z}
\newcommand{\Q}{\mathbb Q}
\newcommand{\C}{\mathbb C}
\newcommand{\D}{\mathbb D}
\newcommand{\F}{\mathbb F}
\newcommand{\sref}[2]{\hyperlink{source:#1:#2}{[S#2]}}
\newcommand{\cataloguescope}[1]{#1}
\begin{document}
% UnsolvedMath ID 7200021; source catalogue code AMR-071-0021
\setcounter{section}{7200020}
\section{Zariski multiplicity conjecture on the topological equisingularity and equimultiplicity of varieties at singular points}\label{7200021}
{\small UnsolvedMath ID 7200021 · Algebraic Geometry}
\subsection{Definitions and mathematical statement}

\paragraph{Shared notation.}$\mathbb N$, $\mathbb Z$, $\mathbb Q$, $\mathbb R$, and $\mathbb C$ denote the natural numbers, integers, rationals, reals, and complex numbers, respectively. All additional parameters and hypotheses are specified in the question.


\paragraph{Statement recorded in the supplied source.}
If two complex hypersurface germs have the same embedded topological singularity type, must their multiplicities at the singular point agree? The multiplicity is the least degree of a nonzero term in a local defining equation. The related family version asks whether topological equisingularity forces equimultiplicity.
\par\smallskip\textit{Formulation references:} \sref{7200021}{1}, \sref{7200021}{2}.
\subsection{Short English statement}
If two complex hypersurface germs have the same embedded topological singularity type, must their multiplicities at the singular point agree? The multiplicity is the least degree of a nonzero term in a local defining equation. The related family version asks whether topological equisingularity forces equimultiplicity.
\subsection{Sources}
\begin{itemize}
\item[\hypertarget{source:7200021:1}{S1}] ULAM.AI and the UnsolvedMath Contributors, UnsolvedMath: A Curated Collection of Open Mathematics Problems, record 7200021 (AMR-071-0021). Catalogue attribution under CC BY 4.0.
\url{https://huggingface.co/datasets/ulamai/UnsolvedMath}.
\item[\hypertarget{source:7200021:2}{S2}] Original problem formulation and mathematical context.
\url{https://en.wikipedia.org/w/index.php?title=List_of_unsolved_problems_in_mathematics&oldid=1366636767#Geometry}.
\end{itemize}
\label{end:7200021}
\end{document}
